% Generated by scripts/build_manuscript.py from sections/speech_forecasting.tex.
% Edit the source, then regenerate; all outcomes are retained.
\section{Speech-energy forecasting with group-level calibration}
\label{sec:speech}

This second real-data domain tests whether a larger number of documented
recording groups permits useful calibration. We use the official Mini Speech
Commands subset \citep{warden2018speech,tensorflow2026mini}, with an eligibility
rule, group split, observable, and model grid specified
before downloading or inspecting the audio.

\paragraph{Observable and sampling units.}
Of 8000 source clips, 7178 have the prescribed mono, 16 kHz, signed 16-bit,
exactly one-second format. They belong to 1630 filename identifiers, using the
source convention before \texttt{\_nohash\_}. All clips from an identifier
remain together. A fixed random permutation assigns 244 identifiers to
training, 163 to validation, 978 to calibration, and 245 to testing. These
contain 1159, 677, 4256, and 1086 clips, respectively. Identifiers denote
recording sessions for web volunteers; they are not biometric confirmation
of distinct people. Confidence statements assume iid group-record bundles
and independence from development data. Identifier separation supports the
evaluation design but does not establish that population sampling assumption.

Each clip is divided into 100 nonoverlapping frames of 160 samples. For frame
$t\in\{0,\ldots,99\}$ and sample index $j\in\{1,\ldots,160\}$,
let $x_{tj}$ be the signed pulse-code modulation (PCM) amplitude divided by
32768. The observed energy is
\[
 e_t=\frac{10\log_{10}\!\left(\max\left\{160^{-1}\sum_{j=1}^{160}
 x_{tj}^2,10^{-6}\right\}\right)+60}{60}\in[0,1].
\]
The implementation clips numerical output to $[0,1]$. A frame is observed at
its end, so this transformation uses no future samples. It represents power
in decibels relative to full scale, floored at $-60$ dB. Multiplying normalized
RMSE by 60 gives RMSE in dB for this floored observable, not waveform error
or error in unfloored acoustic power. The source selects and quality-screens
one-second recordings; our complete-clip eligibility further restricts the
population. This is retrospective forecasting within released clips, without
per-utterance normalization by future means, peaks, or extrema.

Histories contain 4 or 12 frames, horizons are 1, 4, or 12 frames
(10, 40, or 120 ms), and all candidates share origins $11,15,\ldots,87$.
Losses average over origins and clips within each identifier, then weight
identifiers equally. Thousands of clips and windows do not replace the
978 independent calibration units assumed by the theorem.

\paragraph{Development and calibration.}
Three fitting seeds use the same group split. Each samples 4096 training
histories with equal identifier probability and evaluates all validation
groups with equal weights. The prespecified study compares 486 candidates per
fit across eight
families (Table~\ref{tab:speech}). The geometric families share prefix sizes
$2,6,14,30$, standardization or raw inputs, budgets $1,4$, and ridge parameters
$10^{-4},10^{-2}$. RBF and affine affinity multipliers are $0.25,1,4$;
KAHM uses its inherited reconstruction. Conventional regression heads use
ridge regularization without a coefficient-norm constraint; thus their
comparison with the geometric heads does not control statistical capacity.
Appendix~\ref{app:speech-protocol} specifies their complete grids. Every family
has both histories. Selection
minimizes validation MSE across horizons; all 24 selected models are frozen
before held-out losses are evaluated.

For each selected model and horizon, calibration gives one bounded loss in
$[0,1]$ per identifier. The bounded-mean relative-entropy upper bound in
Appendix~\ref{main-app:kl-calibration} uses tail probability $0.05/72$, covering
all $3\times8\times3$ statements jointly. Prespecified tolerances of
6, 12, and 18 dB are explicit resolution levels, not claimed application
requirements. Model choices and calibration formulas remain fixed during
held-out evaluation. Appendix~\ref{app:speech-protocol} describes the scoring
order and its verification.

\begin{table}[t]
\centering\small
\setlength{\tabcolsep}{3pt}
\begin{tabular}{lrrrrrrr}
\toprule
& \multicolumn{3}{c}{Test RMSE in dB: mean (SD)} & \multicolumn{3}{c}{Maximum upper RMSE in dB} & All $\leq18$ dB\\
Family & 10 ms & 40 ms & 120 ms & 10 ms & 40 ms & 120 ms & Fits / 3\\
\midrule
KAHM blocks & 4.00 (0.15) & 7.58 (0.14) & 12.73 (0.06) & 8.15 & 11.64 & 16.37 & 3/3\\
RBF blocks & 3.79 (0.01) & 7.45 (0.01) & 12.72 (0.00) & 7.83 & 11.41 & 16.30 & 3/3\\
Affine-hull blocks & 4.28 (0.43) & 7.77 (0.24) & 12.81 (0.05) & 8.80 & 11.98 & 16.41 & 3/3\\
Ridge AR & 3.49 (0.00) & 7.22 (0.01) & 12.51 (0.01) & 7.53 & 11.19 & 16.15 & 3/3\\
Quadratic & 3.52 (0.01) & 7.24 (0.03) & 12.47 (0.03) & 7.55 & 11.21 & 16.12 & 3/3\\
Gaussian kernel & 3.51 (0.00) & 7.17 (0.01) & 12.39 (0.02) & 7.57 & 11.15 & 16.02 & 3/3\\
Persistence & 3.56 (0.00) & 7.71 (0.00) & 14.70 (0.00) & 7.58 & 11.61 & 18.08 & 0/3\\
Constant & 15.67 (0.01) & 15.63 (0.01) & 15.49 (0.01) & 19.21 & 19.17 & 19.04 & 0/3\\
\bottomrule
\end{tabular}
\caption{Speech-energy forecasting on 245 held-out identifiers. Test means and sample standard deviations use three fits on the same group split. Upper columns take the maximum over those fits, with simultaneous coverage of all 72 selected family--fit--horizon statements under the stated iid group-record assumption. All six learned families pass the 18 dB requirement at every horizon in each fit; none passes the all-horizon 6 or 12 dB requirement. Errors concern the floored energy observable.}
\label{tab:speech}
\end{table}


\paragraph{Findings.}
At 120 ms, mean test RMSE is 12.73 dB for KAHM, 12.51 for ridge AR,
12.47 for quadratic regression, and 12.39 for Gaussian-kernel regression.
The latter three also outperform KAHM at shorter horizons. KAHM improves
on long-horizon persistence, whose mean is 14.70 dB. These are descriptive
accuracy comparisons; no separately calibrated superiority claim is made.
The displayed standard deviations describe fitting variation on the same
245 test groups.

All six learned families satisfy the all-horizon 18 dB upper criterion in
all three fits. Persistence and the constant model do not; no family passes
the all-horizon 6 or 12 dB criterion. Thus the larger calibration cohort can
resolve a declared tolerance under the stated sampling model, while leaving
tighter requirements unresolved. The experiment supports the certification
procedure without establishing that its tolerance is sufficient for an audio
application or that KAHM is the best forecaster. Shared geometry costs and inference times
are archived separately and imply no general computational advantage.

\paragraph{Input scaling on the same recording cohort.}
A descriptive comparison evaluates standardization, raw energy histories,
training min--max scaling to $[-1,1]$, and its clipped variant for all three
geometric families. It retains the same observations, 4096 sampled training
histories per seed, and decoder grids. Each transformation is fitted before
the 1600-history geometry subsample. There are 2688 candidate fits in total;
validation selects one model within each family, seed, and preprocessing arm.
All 36 selections are frozen before test scoring. Because this analysis
reuses the previously examined recording cohort, its results provide no
independent confirmation or new risk certificate. Calibration groups are
not scored in this comparison.

% Generated by experiments/speech_scaling/make_assets.py from retained group losses.
\begin{table}[t]
\centering
\small
\setlength{\tabcolsep}{3.5pt}
\begin{tabular}{lrrrrrr}
\toprule
& Standard & Raw & Min--max & Clipped & $\Delta\%$ & Fits \\
\midrule
KAHM & 22.570 (0.197) & 21.800 (0.441) & 22.251 (0.092) & 22.251 (0.092) & $-1.41$ & 3/3 \\
RBF & 21.466 (0.034) & 21.450 (0.023) & 21.304 (0.132) & 21.304 (0.132) & $-0.76$ & 3/3 \\
Affine hull & 22.493 (0.665) & 22.995 (0.661) & 22.547 (0.104) & 22.547 (0.104) & $+0.24$ & 1/3 \\
\bottomrule
\end{tabular}
\caption{Descriptive speech scaling comparison on the same 245 test groups.
Entries are all-horizon MSE means (sample standard deviations across three fits),
multiplied by $10^3$. Each fit averages the three horizon MSEs with equal weight.
All model choices within an arm use validation data. $\Delta\%$ compares the
min--max and standardized means; negative values favor min--max. ``Fits'' counts
the paired fits with lower min--max MSE. The clipped arm shares training extrema
and clips held-out inputs to $[-1,1]$. This reused-cohort follow-up supplies no
new confidence or certification claim.}
\label{tab:speech-scaling}
\end{table}


Table~\ref{tab:speech-scaling} shows that min--max scaling lowers KAHM's mean
all-horizon MSE by $1.41\%$ relative to standardization, with lower MSE in
each of the three fits. Raw energy histories nevertheless give lower KAHM
MSE than either min--max variant in every fit. RBF min--max MSE decreases by
$0.76\%$ with improvement in all three fits; affine-hull mean MSE increases
by $0.24\%$, with improvement in only one fit. Clipped and unclipped results
agree at the displayed precision. The ridge AR, quadratic, and
Gaussian-kernel regressors from Table~\ref{tab:speech} retain lower test MSE
than every KAHM scaling arm at every horizon and seed. These comparisons
describe the specified implementations and the same 245 test groups.
